% Arquitectura del Sistema SGM-Payroll - Versión Mejorada
% Basada en el contexto actual del sistema desarrollado

\section{Arquitectura del Sistema}

El sistema SGM-Payroll ha sido desarrollado con una arquitectura híbrida multi-capas, diseñada específicamente para manejar la complejidad inherente del procesamiento de nóminas empresariales. La solución combina principios de microservicios, procesamiento asíncrono por fases y almacenamiento estratificado, desplegada en infraestructura containerizada en Uruguay con acceso VPN restringido por MAC.

\subsection{Principios Arquitectónicos}

La arquitectura del sistema se fundamenta en cuatro principios clave:

\begin{itemize}
    \item \textbf{Procesamiento por Fases}: El flujo de nómina se descompone en 5 fases secuenciales, cada una con responsabilidades específicas y almacenamiento optimizado.
    \item \textbf{Almacenamiento Estratificado}: Utilización de PostgreSQL para persistencia, Redis para cache temporal por fases, y almacenamiento en memoria para procesamiento activo.
    \item \textbf{Escalabilidad Horizontal}: Arquitectura preparada para múltiples trabajadores Celery y particionamiento de datos por cliente.
    \item \textbf{Tolerancia a Fallos}: Sistema de reintentos automáticos, rollback por fases y recuperación de estado desde cache.
\end{itemize}

\subsection{Componentes del Sistema}

\subsubsection{Capa de Presentación}
\begin{itemize}
    \item \textbf{React + Vite}: SPA optimizada con componentes modulares por rol de usuario (Analista, Supervisor, Gerente)
    \item \textbf{Tailwind CSS}: Sistema de diseño consistente con componentes reutilizables
    \item \textbf{Streamlit Dashboard}: Visualizaciones interactivas para KPIs y métricas de performance
\end{itemize}

\subsubsection{Capa de API y Lógica de Negocio}
\begin{itemize}
    \item \textbf{Django REST Framework}: API RESTful con ViewSets personalizados, serializers optimizados y filtrado por cliente
    \item \textbf{Sistema de Autenticación}: Control de acceso basado en roles con asignación automática de analistas responsables
    \item \textbf{Validadores Modulares}: Validación de RUT, formatos de archivo y consistencia de datos por fase
\end{itemize}

\subsubsection{Capa de Procesamiento Asíncrono}
\begin{itemize}
    \item \textbf{Celery Workers}: Procesamiento distribuido con chains de tareas por fase
    \item \textbf{Redis Broker}: Orquestación de tareas con patrones de retry y circuit breaker
    \item \textbf{Flower Monitoring}: Supervisión en tiempo real de workers y cola de tareas
\end{itemize}

\subsubsection{Capa de Almacenamiento}
\begin{itemize}
    \item \textbf{PostgreSQL}: Base de datos principal con modelos optimizados para consultas por cliente y período
    \item \textbf{Redis DB Especializado}: Cache estructurado por fases con TTL diferenciado según criticidad de datos
    \item \textbf{Sistema de Logs}: Logging estructurado con rotación automática y políticas de retención
\end{itemize}

\subsection{Arquitectura de Procesamiento por Fases}

El sistema implementa un flujo de procesamiento innovador basado en 5 fases secuenciales, cada una optimizada para minimizar el uso de memoria y maximizar la eficiencia:

\subsubsection{Fase 1: Análisis de Headers y Mapeo}
\begin{itemize}
    \item \textbf{Objetivo}: Extraer y mapear conceptos del Libro de Remuneraciones Talana
    \item \textbf{Almacenamiento Redis}: Solo metadatos de headers y mapeos (≈5KB por archivo)
    \item \textbf{Optimización}: No se almacenan datos de empleados, reduciendo uso de memoria en 85\%
    \item \textbf{Estructura}: \texttt{sgm:nomina:\{cliente\}:\{cierre\}:fase1\_headers:\{tipo\}}
\end{itemize}

\subsubsection{Fase 2: Movimientos Talana}
\begin{itemize}
    \item \textbf{Objetivo}: Procesar ausentismos, ingresos y finiquitos
    \item \textbf{Filtrado Inteligente}: Solo empleados existentes en Libro de Remuneraciones
    \item \textbf{Almacenamiento}: Movimientos filtrados en PostgreSQL, cache selectivo en Redis
\end{itemize}

\subsubsection{Fase 3-4: Archivos del Analista}
\begin{itemize}
    \item \textbf{Validación Cruzada}: Empleados deben existir en lista del libro
    \item \textbf{Mapeo de Headers}: Sincronización de términos entre analista y Talana
    \item \textbf{Escalabilidad}: Filtrado automático reduce datos innecesarios
\end{itemize}

\subsubsection{Fase 5: Verificación y Discrepancias}
\begin{itemize}
    \item \textbf{Comparación de Datos}: Talana vs Analista en movimientos y conceptos-valor
    \item \textbf{Cache Temporal}: Solo durante verificación activa, limpieza automática post-proceso
    \item \textbf{Persistencia}: Discrepancias finales en PostgreSQL para auditoría
\end{itemize}

\subsection{Gestión de Cache Estratificado}

El sistema implementa una estrategia de cache híbrida con Redis DB2 dedicado exclusivamente a nómina:

\begin{verbatim}
Estructura Organizada:
sgm:nomina:{cliente_id}:{cierre_id}:{componente}:{tipo}

Componentes por TTL:
├── fase1_headers:* (TTL: 24h) - Mapeos críticos
├── estado:procesamiento (TTL: 2h) - Estado temporal  
├── archivo:libro_excel (TTL: 2h) - Datos temporales
└── fase5_verificacion:* (TTL: 1h) - Cache volátil
\end{verbatim}

\subsubsection{Ventajas del Cache Estratificado}
\begin{itemize}
    \item \textbf{Uso Eficiente de Memoria}: Escalado controlado por fases vs crecimiento exponencial
    \item \textbf{TTL Diferenciado}: Retención optimizada según criticidad de datos
    \item \textbf{Limpieza Automatizada}: Cleanup por patrones y fases completadas
    \item \textbf{Debugging Mejorado}: Keys autodescriptivas con estructura consistente
\end{itemize}

\subsection{Patrón de Integración API-Cache-BD}

El sistema implementa un patrón de 3 capas para optimizar performance y consistencia:

\begin{enumerate}
    \item \textbf{API Layer}: Validación de entrada y autorización por cliente
    \item \textbf{Cache Layer}: Redis como coordinador de flujo entre fases
    \item \textbf{Persistence Layer}: PostgreSQL para almacenamiento definitivo y auditoría
\end{enumerate}

\subsubsection{Flujo de Datos Optimizado}
\begin{verbatim}
Upload Excel → Headers Redis + Empleados PostgreSQL
Mapeo Manual → Cache selectivo conceptos en uso
Confirmación → Cleanup Redis + Actualización PostgreSQL
\end{verbatim}

\subsection{Escalabilidad y Performance}

\subsubsection{Métricas de Optimización Logradas}
\begin{itemize}
    \item \textbf{Reducción de Memoria Redis}: 85\% menos uso vs almacenamiento completo
    \item \textbf{Filtrado Automático}: Solo datos relevantes por fase (empleados existentes)
    \item \textbf{Procesamiento Incremental}: 133 empleados × 55 conceptos procesados eficientemente
    \item \textbf{Cache Selectivo}: TTL diferenciado según criticidad (1h-24h)
\end{itemize}

\subsubsection{Capacidad de Escalado}
\begin{itemize}
    \item \textbf{Horizontal}: Workers Celery distribuidos por cliente
    \item \textbf{Vertical}: Cache Redis particionado por DB (DB2 dedicado nómina)
    \item \textbf{Temporal}: Cleanup automático libera recursos post-proceso
\end{itemize}

\subsection{Seguridad y Acceso}

\subsubsection{Infraestructura}
\begin{itemize}
    \item \textbf{Despliegue}: Docker Swarm en servidor Uruguay
    \item \textbf{Acceso}: VPN obligatorio + whitelist MAC addresses
    \item \textbf{Datos}: Anonimización automática para testing
\end{itemize}

\subsubsection{Aplicación}
\begin{itemize}
    \item \textbf{Autenticación}: Django auth con roles diferenciados
    \item \textbf{Autorización}: Filtrado automático por cliente asignado
    \item \textbf{Auditoría}: Logs estructurados de todas las operaciones
\end{itemize}

Esta arquitectura híbrida por fases representa una innovación específica para el dominio de nóminas, donde la optimización del uso de recursos y la organización por flujo de negocio son críticas para la escalabilidad del sistema.
